\chapter{Quick Reference}
\label{ch:cheatsheet}

This appendix provides a compact reference for Novus syntax and common patterns.

\section{Variables and Constants}

\begin{lstlisting}
// ... (variable examples - illustrative)
var x = 10
x = 20

// Immutable binding
let y = 42

// Compile-time constant
pub const MAX_SIZE: u32 = 1024

// Type annotation (when needed)
let z: i32 = 100
var buffer = [0u8; 256]   // Size inferred
\end{lstlisting}

\section{Primitive Types}

\begin{center}
\begin{tabular}{lll}
\textbf{Type} & \textbf{Size} & \textbf{Range/Description} \\
\hline
\texttt{i8} & 8-bit & $-128$ to $127$ \\
\texttt{i16} & 16-bit & $-32{,}768$ to $32{,}767$ \\
\texttt{i32} & 32-bit & $\pm 2$ billion (default) \\
\texttt{i64} & 64-bit & Very large signed \\
\texttt{u8} & 8-bit & $0$ to $255$ \\
\texttt{u16} & 16-bit & $0$ to $65{,}535$ \\
\texttt{u32} & 32-bit & $0$ to $4$ billion \\
\texttt{u64} & 64-bit & Very large unsigned \\
\texttt{f32} & 32-bit & Single-precision float \\
\texttt{f64} & 64-bit & Double-precision float \\
\texttt{fixed16} & 16-bit & 8.8 fixed-point \\
\texttt{fixed32} & 32-bit & 16.16 fixed-point \\
\texttt{bool} & 8-bit & \texttt{true} or \texttt{false} \\
\texttt{*T} & 32-bit & Raw pointer (nullable) \\
\texttt{\&T} & 32-bit & Reference (non-null) \\
\texttt{\&var T} & 32-bit & Mutable reference \\
\end{tabular}
\end{center}

\section{Literals}

\begin{lstlisting}
// ... (literal examples - illustrative)
42          // i32 (default)
42u32       // u32
42i8        // i8
0xFF        // Hexadecimal (C-style)
$FF         // Hexadecimal (Amiga-style)
%1010       // Binary
1_000_000   // Underscores for readability

// Floats
3.14        // f64 (default)
3.14f32     // f32

// Strings
"Hello"              // String literal
"Line1\nLine2"       // With escape
f"Value: {x}"        // F-string (interpolation)

// Characters
'A'         // u8 (65)
'\n'        // Newline
'\x1B'      // Hex byte (ESC)

// Other
true, false // Booleans
null        // Null pointer
\end{lstlisting}

\section{Operators}

\begin{center}
\begin{tabular}{ll|ll}
\textbf{Op} & \textbf{Meaning} & \textbf{Op} & \textbf{Meaning} \\
\hline
\texttt{+} & Add & \texttt{\&\&} & Logical AND \\
\texttt{-} & Subtract & \texttt{||} & Logical OR \\
\texttt{*} & Multiply & \texttt{!} & Logical NOT \\
\texttt{/} & Divide & \texttt{\&} & Bitwise AND \\
\texttt{\%} & Modulo & \texttt{|} & Bitwise OR \\
\texttt{==} & Equal & \texttt{\^{}} & Bitwise XOR \\
\texttt{!=} & Not equal & \texttt{\~{}} & Bitwise NOT \\
\texttt{<} & Less than & \texttt{<<} & Left shift \\
\texttt{>} & Greater than & \texttt{>>} & Right shift \\
\texttt{<=} & Less or equal & \texttt{..} & Range (exclusive) \\
\texttt{>=} & Greater or equal & \texttt{..=} & Range (inclusive) \\
\end{tabular}
\end{center}

Compound: \texttt{+=}, \texttt{-=}, \texttt{*=}, \texttt{/=}, \texttt{\%=}, \texttt{\&=}, \texttt{|=}, \texttt{\^{}=}, \texttt{<<=}, \texttt{>>=}

Increment/Decrement: \texttt{++x}, \texttt{x++}, \texttt{--x}, \texttt{x--}

\section{Control Flow}

\begin{lstlisting}
// If-else
if condition {
    // ...
} else if other {
    // ...
} else {
    // ...
}

// If as expression
let max = if a > b { a } else { b }

// Match
match value {
    1 => handle_one(),
    2 | 3 => handle_two_or_three(),
    4..=10 => handle_range(),
    _ => handle_default(),
}

// While loop
while condition {
    // ...
}

// For loop
for i in 0..10 {
    // i goes 0 to 9
}

for item in collection {
    // iterate items
}

// Loop (infinite)
loop {
    if done { break }
}

// Break and continue
break           // Exit loop
continue        // Next iteration
break value     // Return value from loop
\end{lstlisting}

\section{Functions}

\begin{lstlisting}
// Basic function
fn add(a: i32, b: i32) -> i32 {
    return a + b
}

// No return value
fn print_hello() {
    println("Hello")
}

// Public function
pub fn api_function() -> i32 {
    return 42
}

// Generic function
fn swap<T>(a: &var T, b: &var T) {
    let tmp = *a
    *a = *b
    *b = tmp
}

// With trait bounds
fn print_all<T>(items: &[T]) where T: Display {
    for item in items {
        println(item.to_string())
    }
}
\end{lstlisting}

\section{Structs and Enums}

\begin{lstlisting}
// Struct
struct Point {
    x: i32,
    y: i32,
}

// Create instance
let p = Point { x: 10, y: 20 }

// Impl block
impl Point {
    pub fn new(x: i32, y: i32) -> Point {
        return Point { x: x, y: y }
    }

    pub fn distance(&self, other: &Point) -> f32 {
        // ...
    }
}

// Enum
enum Direction {
    Up,
    Down,
    Left,
    Right,
}

// Enum with data
enum Message {
    Quit,
    Move(i32, i32),
    Write(String),
}

// Match enum
match msg {
    Message::Quit => exit(),
    Message::Move(x, y) => move_to(x, y),
    Message::Write(s) => println(s),
}
\end{lstlisting}

\section{Option and Result}

\begin{lstlisting}
// ... (Option/Result patterns - illustrative)
let maybe: Option<i32> = Option::Some(42)
let none: Option<i32> = Option::None

match maybe {
    Option::Some(x) => use_value(x),
    Option::None => handle_missing(),
}

let value = maybe.unwrap_or(0)

// Result
let result: Result<i32, Error> = Result::Ok(42)
let error: Result<i32, Error> = Result::Err(Error::NotFound)

match result {
    Result::Ok(x) => use_value(x),
    Result::Err(e) => handle_error(e),
}

// ? operator (propagate error)
fn process() -> Result<i32, Error> {
    let file = open(path)?   // Returns Err early
    let data = read(&file)?
    Result::Ok(parse(data))
}
\end{lstlisting}

\section{Memory and References}

\begin{lstlisting}
// Immutable reference
let r: &i32 = &x

// Mutable reference
let r: &var i32 = &var x
*r = 10  // Modify through reference

// Raw pointer
let p: *u8 = (*u8)address

// Dereference
let value = *r
let byte = *p

// Null check
if ptr != null {
    // safe to use
}

// defer (cleanup)
fn process() {
    let resource = acquire()
    defer { release(resource) }
    // ... use resource
}  // release() called here

// Drop trait (RAII)
impl Drop for MyHandle {
    fn drop(&var self) {
        // Automatic cleanup
    }
}
\end{lstlisting}

\section{RAII Patterns}

\begin{center}
\begin{tabular}{lp{6.5cm}}
\textbf{Pattern} & \textbf{Use} \\
\hline
\texttt{Handle} & Wrap OS resources (\texttt{WindowHandle}, \texttt{FileHandle}) \\
\texttt{Guard} & Temporary locks (\texttt{SemaphoreGuard}, \texttt{BlitterGuard}) \\
\texttt{Builder} & Accumulate config (\texttt{MUIAppBuilder}) \\
\end{tabular}
\end{center}

\begin{lstlisting}
// Handle: wrap resource, auto-cleanup
var window = WindowBuilder::new().build()?
// window closed automatically on drop

// Guard: temporary exclusive access
var lock = semaphore.obtain()
// ReleaseSemaphore() called on drop

// Builder: configure then create
var app = MUIAppBuilder::new()
    .title("App")
    .build()?
\end{lstlisting}

\section{Arrays and Collections}

\begin{lstlisting}
// Fixed array (size inferred)
let arr = [1, 2, 3]
let zeros = [0; 100]  // 100 zeros
let first = arr[0u32]

// Vec (dynamic array)
var v: Vec<i32> = Vec::new()
v.push(1)
v.push(2)
let len = v.len()
let item = v.get(0u32)  // Option<T>

// Tuple
let point: (i32, i32) = (10, 20)
let (x, y) = point  // Destructure
\end{lstlisting}

\section{Modules and Imports}

\begin{lstlisting}
// ... (import examples - illustrative)
from std::collections::vec import Vec
from std::io::file import println, read_file

// Multiple imports
from std::core import Option, Result, Drop

// Use in code
let v: Vec<i32> = Vec::new()
\end{lstlisting}

\section{Traits}

\begin{lstlisting}
trait Printable {
    fn print(&self)
}
\end{lstlisting}

Implementation:

\begin{lstlisting}
// ... (with Point and Printable defined above)
impl Printable for Point {
    fn print(&self) {
        println(f"({self.x}, {self.y})")
    }
}
\end{lstlisting}

Common traits: \texttt{Drop} (destructor), \texttt{Clone}, \texttt{Eq}, \texttt{Hash}.

\section{Common Attributes}

\begin{lstlisting}[language=bash]
@test                     // Test function
@test(skip = "reason")    // Skipped test
@benchmark                // Benchmark function
@inline                   // Hint: inline this function
@noinline                 // Prevent inlining
@packed                   // No struct padding
@export                   // Keep function, don't mangle name
@atomic                   // Wrap in Forbid/Permit
#[derive(Eq, Hash)]       // Auto-implement traits
#[stack_size(65536)]      // Set program stack size
\end{lstlisting}

\section{Intrinsics}

\begin{lstlisting}[language=bash]
@sizeof(T)      // Size of type in bytes
@zeroed(T)      // Zero-initialized value
\end{lstlisting}

Note: \texttt{alignof()} is only available in inline assembly blocks.

\section{Compiler Commands}

\begin{lstlisting}[language=bash]
# Compile single file
novus compile file.novus -o output

# Build project
novus build

# Run tests
novus test file.novus

# Format code
novus fmt

# With options
novus compile file.novus -o out --release --cpu 68040
\end{lstlisting}

\section{C to Novus Quick Reference}

\begin{center}
\begin{tabular}{ll}
\textbf{C} & \textbf{Novus} \\
\hline
\texttt{int x = 10;} & \texttt{var x = 10} \\
\texttt{const int X = 10;} & \texttt{const X: i32 = 10} \\
\texttt{if (ptr)} & \texttt{if ptr != null} \\
\texttt{switch/case/break} & \texttt{match \{ pat => \}} \\
\texttt{for(i=0;i<n;i++)} & \texttt{for i in 0..n} \\
\texttt{NULL} & \texttt{null} \\
\texttt{struct S \{...\};} & \texttt{struct S \{...\}} \\
\texttt{typedef struct} & (not needed) \\
\texttt{ptr->field} & \texttt{ptr.field} (auto-deref) \\
\texttt{malloc/free} & Use \texttt{Vec}, \texttt{MemoryBlock}, or allocator \\
\texttt{sizeof(T)} & \texttt{@sizeof(T)} \\
\texttt{errno} & \texttt{Result::Err(DosError::...)} \\
\end{tabular}
\end{center}

\section{Escape Sequences}

\begin{center}
\begin{tabular}{ll}
\textbf{Sequence} & \textbf{Meaning} \\
\hline
\texttt{\textbackslash n} & Newline \\
\texttt{\textbackslash r} & Carriage return \\
\texttt{\textbackslash t} & Tab \\
\texttt{\textbackslash 0} & Null byte \\
\texttt{\textbackslash\textbackslash} & Backslash \\
\texttt{\textbackslash "} & Double quote \\
\texttt{\textbackslash '} & Single quote \\
\texttt{\textbackslash xNN} & Hex byte \\
\end{tabular}
\end{center}
